% lecture21.tex — Integration on Manifolds

\section{Integration on Manifolds}

\subsection{Change of Variables in $\R^m$}

Recall from vector calculus that, if $f \colon V \to U$ is a diffeomorphism of open sets in $\R^m$ and $a$ is an integrable function on $U$, then
\[
    \int_U a\, dx_1 \cdots dx_m = \int_V (a \circ f)\, |\!\det(df)|\, dy_1 \cdots dy_m.
\]
Notice the similarity with the pullback formula:
\[
    f^*(a\, dx_1 \wedge \cdots \wedge dx_m) = (a \circ f)\, \det(df)\, dy_1 \wedge \cdots \wedge dy_m.
\]
The two differ only in the sign $\det(df)$ vs.\ $|\det(df)|$. It follows that:

\begin{theorem}{Change of Variables in $\R^m$}{change-of-variables-Rm}
    If $f \colon V \to U$ is an orientation-preserving diffeomorphism of open sets in $\R^m$ (or $H^m$) and $\omega$ is an integrable $m$-form on $U$, then
    \[
        \int_U \omega = \int_V f^*\omega.
    \]
    If $f$ reverses the orientation, then $\int_U \omega = -\int_V f^*\omega$.
\end{theorem}

\subsection{Integral of an $m$-Form over a Manifold}

\begin{definition}{Integral of an $m$-Form}{integral-mform}
    Let $M$ be an oriented $m$-manifold and let $\omega$ be a smooth, compactly supported $m$-form on $M$. Define the \emph{integral of $\omega$ over $M$} to be
    \[
        \int_M \omega := \sum_i \int_{U_i} \rho_i \omega,
    \]
    a finite sum, where $\{\rho_i\}$ is a partition of unity subordinate to parametrizable open subsets and each integral is defined in some local parametrization.
\end{definition}

The change of variables formula ensures that the integral is well-defined, independent of the choice of parametrizations. In particular:

\begin{theorem}{Naturality of the Integral}{integral-naturality}
    If $f \colon N \to M$ is an orientation-preserving diffeomorphism, then
    \[
        \int_M \omega = \int_N f^*\omega
    \]
    for any compactly supported smooth $m$-form $\omega$ (where $m = \dim M = \dim N$).
\end{theorem}

That is, the integral transforms naturally!

\subsection{Integration over Submanifolds}

Just like we can integrate an $m$-form over $M$, we can integrate any $p$-form $\omega$ over an oriented $p$-dimensional submanifold $i \colon P \hookrightarrow M$ by restricting $\omega$ to $P$:
\[
    \int_P \omega := \int_P i^*\omega.
\]

\begin{example}{Line Integral of a 1-Form on $\R^3$}{line-integral}
    Suppose $\omega = f_1\, dx_1 + f_2\, dx_2 + f_3\, dx_3$ is a smooth 1-form on $\R^3$, and $\gamma \colon I \to \R^3$ is a diffeomorphism of $I = [0,1]$ onto a curve $C = \gamma(I)$. Then
    \[
        \int_C \omega = \int_I \gamma^* \omega = \int_0^1 \sum_{i=1}^{3} f_i(\gamma(t))\, d\gamma_i = \int_0^1 \sum_{i=1}^{3} f_i(\gamma(t))\, \frac{d\gamma_i}{dt}(t)\, dt.
    \]
\end{example}

\begin{example}{Surface Integral of a 2-Form on $\R^3$}{surface-integral}
    Suppose
    \[
        \omega = f_1\, dx_2 \wedge dx_3 + f_2\, dx_3 \wedge dx_1 + f_3\, dx_1 \wedge dx_2
    \]
    is a compactly supported smooth 2-form on $\R^3$. Let $S \subset \R^3$ be a surface, which we assume for simplicity to be the graph of a function $G \colon \R^2 \to \R$, i.e.\
    \[
        S = \{(x_1, x_2, x_3) \in \R^3 \mid x_3 = G(x_1, x_2)\}.
    \]
    We use the parametrization
    \[
        h \colon \R^2 \to S, \qquad (x_1, x_2) \mapsto (x_1, x_2, G(x_1, x_2)).
    \]
    Then $h^* dx_3 = dG = \tfrac{\partial G}{\partial x_1} dx_1 + \tfrac{\partial G}{\partial x_2} dx_2$, so
    \begin{align*}
        \int_S \omega &= \int_{\R^2} h^*\omega \\
        &= \int_{\R^2} (f_1 \circ h)\, dx_2 \wedge dG + (f_2 \circ h)\, dG \wedge dx_1 + (f_3 \circ h)\, dx_1 \wedge dx_2 \\
        &= \int_{\R^2} \left((f_1 \circ h)\!\left(-\tfrac{\partial G}{\partial x_1}\right) + (f_2 \circ h)\!\left(-\tfrac{\partial G}{\partial x_2}\right) + (f_3 \circ h)\right) dx_1 \wedge dx_2.
    \end{align*}
    Note that
    \[
        \vec{n} := \left(-\tfrac{\partial G}{\partial x_1}, -\tfrac{\partial G}{\partial x_2}, 1\right)
    \]
    is normal to $T_x S = \mathrm{Span}\{(1, 0, \tfrac{\partial G}{\partial x_1}), (0, 1, \tfrac{\partial G}{\partial x_2})\}$ at every $x \in S$. It follows that we can rewrite the integral as
    \[
        \int_S \omega = \int_{\R^2} (\vec{F} \cdot \vec{u})\, dA,
    \]
    where $\vec{F} = (f_1, f_2, f_3)$, $\vec{u} = \vec{n}/|\vec{n}|$, and $dA = |\vec{n}|\, dx_1 \wedge dx_2$ — a form familiar from vector calculus.
\end{example}
